\subsection{Вступ до теоретичного розділу: Геометричні перетворення та масштабування зображень}

Цифровий ворпінг (digital image warping) — це галузь обробки зображень, що динамічно розвивається та спеціалізується на геометричних перетвореннях цифрових даних. Геометричне перетворення визначається як операція, яка перевизначає просторові відношення між точками на зображенні. Хоча цей термін часто викликає асоціації із високодеформованими візуальними ефектами, такі трансформації варіюються від найпростіших операцій (перенесення, обертання або масштабування) до складних перетворень. Оскільки всі ці процеси застосовують геометричні трансформації, у сучасній літературі терміни «ворпінг» та «геометричне перетворення» використовуються як взаємозамінні. 

Історично перші геометричні перетворення виконувалися над неперервними (аналоговими) зображеннями за допомогою оптичних систем. Попри високу швидкість роботи, такі системи були обмежені у гнучкості та можливостях контролю, що зумовило повний перехід сучасних алгоритмів на роботу виключно з цифровими (дискретними) зображеннями. Найперші дослідження у цій сфері виникли в середині 1960-х років у програмах космічних спостережень NASA (таких як Landsat та Skylab) для вирішення задач суміщення зображень (image registration) — вирівнювання знімків однієї ділянки, зроблених у різний час або різними сенсорами, задля компенсації спотворень. Згодом ці методи поширилися на комп'ютерний зір, медичну візуалізацію та комп'ютерну графіку.

Прямим, найважливішим і найбільш поширеним частковим випадком геометричних перетворень є \textbf{масштабування зображень}. Масштабування полягає у представленні та зміні розміру оригінального зображення шляхом використання більшої або меншої кількості пікселей. Головна мета цього процесу — отримати зображення іншого розміру, максимально зберігши оригінальний вміст із мінімальними втратами якості. Масштабування реалізується у двох протилежних напрямках:
\begin{itemize}
	\item \textbf{Зменшення масштабу (downscaling):} процес стиснення, за якого розмір вхідного зображення високої роздільної здатності (HR) зменшується для отримання цільового зображення низької роздільної здатності (LR). Ця операція є фундаментальною для швидкого перегляду або обміну даними.
	\item \textbf{Збільшення масштабу (upscaling):} процес покращення, в якому розмір вхідного зображення низької роздільної здатності (LR) збільшується для відновлення або отримання цільового зображення високої роздільної здатності (HR). Це дозволяє виділити та підкреслити важливі деталі об'єктів.
\end{itemize}

Зміна розміру зображення із гарантованим збереженням якості має величезне прикладне значення у багатьох галузях: від повсякденної діяльності до передових наук, включаючи астрономію, біологію, медицину, дистанційне зондування, геймінг, автономне водіння, аерофотозйомку, моніторинг трафіку та відеоспостереження. Окрім цивільних та наукових завдань, аналіз ефектів масштабування є критично важливим для кібербезпеки та протидії зловмисним атакам, наприклад, для виявлення автоматично згенерованих камуфляжних зображень, чия візуальна семантика радикально змінюється після зміни розміру.

Оскільки цифрове масштабування супроводжується просторовими деформаціями дискретної матриці пікселів, воно неминуче викликає появу візуальних артефактів. Для забезпечення якісного довільного масштабування було розроблено низку різноманітних методів, що покликані запобігати появі подібних спотворень. Постійне вдосконалення швидких та ефективних алгоритмів геометричного перетворення сьогодні стимулюється доступністю високопродуктивного апаратного забезпечення, що відкриває шлях до точної обробки растрових даних надвисокої роздільної здатності.